4. Experimental Design and Evaluation Protocol
4.1 Training and Test Partition
The experimental design follows a strict chronological holdout structure. The full dataset spans January 1, 2010 through December 31, 2014, yielding 43,824 hourly observations prior to feature construction. Following the creation of lagged and rolling-window features with a maximum lookback of 168 hours, the first 168 observations are dropped to avoid partially undefined feature rows, leaving 43,656 complete observations.
The training period covers January 1, 2010 through December 31, 2013 (34,896 observations). The test period covers the full calendar year 2014 (8,760 observations). No 2014 observations are used for model fitting. Any hyperparameter selection procedure must be documented separately and must not use held-out 2014 outcomes. The training leakage count is zero.
Within the test period, earlier observed load values are used as lagged inputs for later test hours where chronologically valid — for example, the load observed at hour $t - 24$ is a legitimate input when predicting hour $t$. This is a standard property of autoregressive feature structures and does not constitute training leakage, as no model parameters are updated using test-period data.
All three trained models — Linear Regression, GBoost, and QR-GBT — are fitted exclusively on training-period observations. The three naive baselines (Persistence-1h, Naive-24h, Naive-168h) require no fitting procedure. The PJM Day-Ahead forecast is obtained as an external benchmark and is not produced or fitted by this study.
4.2 Evaluation Windows
Forecast performance is evaluated across three defined windows, each serving a distinct purpose in the evaluation design.
Full-year 2014 test set. The primary evaluation window is the complete 8,760-hour 2014 test year. Full-year metrics provide an aggregate characterization of model performance across the full range of operating conditions present in 2014, including moderate-load periods, shoulder seasons, summer peaks, and the January cold event. Full-year metrics are reported for all models.
January 6–8, 2014 Polar Vortex cold-event window. The stress-test evaluation window is the 72-hour period from 00:00 EPT January 6, 2014 through 23:00 EPT January 8, 2014. The January 6–8, 2014 Polar Vortex produced a near-annual-peak RTO load of 140,510.2 MW at Jan 7 18:00 EPT, equal to 99.18% of the 2014 annual peak of 141,677.9 MW recorded on Jun 17 17:00 EPT. This window is defined as a near-annual-peak cold-weather stress event and evaluated separately to assess whether full-year aggregate metrics adequately characterize model behavior under extreme cold conditions. Event-window metrics are reported for all models.
June 16–18, 2014 summer near-peak comparison window. The comparison window is the 72-hour period from 00:00 EPT June 16, 2014 through 23:00 EPT June 18, 2014, which contains the 2014 annual peak of 141,677.9 MW. This window is evaluated in parallel with the cold-event window to distinguish between general peak-load forecasting difficulty and the specific challenge of extreme cold-weather generalization. Comparing results across the two near-peak windows — one cold-driven, one heat-driven — allows the evaluation to assess whether any performance degradation observed during the polar vortex window is attributable to peak-load conditions generally or to the out-of-distribution character of the cold event specifically.
Top 1% load hours. As a supplementary high-load subset, the top 1% of 2014 test hours by observed load level (approximately 88 hours) may be used in Section 5 to characterize forecast calibration across the upper tail of the 2014 load distribution. This subset is not a primary evaluation window.
All weather-informed model results across all evaluation windows use same-hour ERA5 retrospective reanalysis weather input. This is an idealized retrospective weather-information setting and is not an operational forecast-weather setting. The evaluation does not quantify degradation from weather forecast error, and results should not be interpreted as achievable under real-time operational conditions.
4.3 Point Forecast Evaluation Metrics
Point forecast accuracy is evaluated using two primary metrics: mean absolute error (MAE) and root mean squared error (RMSE). Both metrics are computed over all hours within each evaluation window.
MAE is the average of the absolute differences between observed and predicted load across all hours in the evaluation window, expressed in megawatts. MAE is the primary point forecast accuracy metric throughout this study.
RMSE is the square root of the average squared differences between observed and predicted load. RMSE applies greater weight to large errors than MAE. Because the polar vortex event produces concentrated large errors at peak hours, RMSE provides additional sensitivity to these tail errors relative to MAE.
Error sign convention. Forecast error at hour $t$ is defined as actual minus predicted: $e_t = y_t - \hat{y}_t$. A positive error indicates underprediction; a negative error indicates overprediction. This convention is applied consistently across all error reporting and directional analysis.
Event-peak point error. For the polar vortex evaluation window, the point forecast error at the peak hour (Jan 7 18:00 EPT, observed load 140,510.2 MW) is reported individually for each model in Table 5. This single-hour error is distinct from the window-level MAE and RMSE and is reported alongside the probabilistic interval inclusion check described in Section 4.4.
Point forecast results for all models and all evaluation windows are reported in Table 3.
4.4 Probabilistic Forecast Evaluation Metrics
Probabilistic forecast evaluation is applied to the QR-GBT model. All probabilistic evaluation uses the post-hoc rearranged quantile estimates, as described in Section 3.5. As noted there, quantile crossing was detected in 5,786 of 8,760 full-year 2014 test hours (66% of test hours) prior to rearrangement. The rearrangement correction is a methodological limitation and does not guarantee well-calibrated intervals.
The following metrics are defined and reported in Section 5.
Median forecast MAE (q50 MAE). The MAE of the QR-GBT q50 (median) forecast, computed using the same definition as the point forecast MAE in Section 4.3. This allows direct comparison of the QR-GBT median forecast with the GBoost and Linear Regression point forecasts on a common accuracy scale.
Mean pinball loss. Mean pinball loss is computed by averaging the pinball loss — defined in Section 3.1 — over all evaluated hours and over the seven estimated quantile levels: q01, q05, q10, q50, q90, q95, and q99. Lower mean pinball loss indicates better overall probabilistic forecast skill across the full quantile set. Mean pinball loss is reported by evaluation window.
Empirical prediction interval coverage. Empirical coverage is the fraction of evaluation hours for which the observed load falls within the estimated interval bounds. For the nominal 90% prediction interval, the bounds are the q05 lower quantile and the q95 upper quantile. For the nominal 98% prediction interval, the bounds are the q01 lower quantile and the q99 upper quantile. A well-calibrated model should achieve empirical coverage close to its nominal level. Systematic undercoverage — empirical coverage materially below the nominal level — indicates that the estimated intervals are too narrow relative to the actual conditional uncertainty. Empirical coverage rates are reported by evaluation window in Section 5.
Winkler score. The Winkler score evaluates prediction interval sharpness and interval violations jointly. It rewards narrower intervals when observations fall inside the interval bounds, and penalizes observations that fall outside the interval in proportion to the exceedance distance. The penalty magnitude scales with the nominal miscoverage rate, so that violations of higher-confidence intervals are penalized more heavily. Lower mean Winkler score indicates better joint sharpness and coverage. The Winkler score is reported for the nominal 90% prediction interval across all evaluation windows.
Event-peak interval inclusion check. For the polar vortex peak hour (Jan 7 18:00 EPT, observed load 140,510.2 MW), the evaluation records whether the observed load falls within each estimated prediction interval. This is a binary inclusion check, reported in Table 5 for both the nominal 90% interval (q05–q95 bounds) and the nominal 98% interval (q01–q99 bounds). The post-rearrangement q95 and q99 estimates at the peak hour are also reported in Table 5 to allow direct comparison with the observed peak load.
Full probabilistic evaluation results for all windows are reported in Table 4. Event-peak probabilistic results are reported in Table 5.
4.5 Benchmark Comparisons and Interpretation Rules
The evaluation compares multiple models across point forecast and probabilistic forecast tasks. Several interpretation rules govern how these comparisons should be read.
Naive baseline comparisons. The three naive baselines — Persistence-1h, Naive-24h, and Naive-168h — require no training and use only their defining lagged load values. These baselines provide reference points for assessing whether trained models add value beyond simple load-memory rules. These comparisons are evaluated using MAE across all windows.
PJM Day-Ahead comparison. The PJM Day-Ahead forecast is an external benchmark produced operationally by PJM using PJM's internal forecasting system and inputs not controlled by this study. It is included to provide operational context for the point forecast results. GBoost and QR-GBT use same-hour ERA5 retrospective reanalysis weather input, which gives these models an informational advantage not available under real-time operational conditions. Accordingly, comparisons with PJM Day-Ahead are retrospective and should not be interpreted as evidence of operational superiority. The purpose of the comparison is to contextualize the magnitude of retrospective gradient boosted model errors relative to those of an operational forecasting system during the same event.
Full-year versus event-window comparison. The comparison between full-year and event-window metrics is used to assess whether aggregate annual evaluation masks stress-window behavior. Section 5 reports both sets of metrics for all applicable models, and the two should be read in conjunction rather than treated as independent summaries.
Retrospective weather caveat. All comparisons involving GBoost or QR-GBT are subject to the constraint that these models use same-hour ERA5 retrospective reanalysis weather input. This is not an operational forecast-weather setting. No comparison in this study quantifies the additional error that would arise from replacing retrospective reanalysis weather with real-time NWP forecast weather.
4.6 Reproducibility and Leakage Controls
All primary data inputs used in this study are publicly available. PJM RTO hourly metered load data are available from PJM Data Miner 2. ERA5 reanalysis fields are available from the Copernicus Climate Data Store [CITATION: ERA5 reanalysis]. The ERA5 great_lakes_core spatial aggregation procedure is defined in src/data/process_era5_2014.py as described in Section 2.2.
The following leakage controls are applied throughout the experimental design:


No 2014 observations are used for model fitting.


Lag and rolling features are constructed using only observations at hours strictly prior to the prediction target hour $t$, respecting chronological ordering at all times.


The 168-hour trim at the start of the dataset ensures that no partially defined feature rows are included in training or evaluation.


The training and test periods are defined by calendar year boundaries and are non-overlapping.


These controls ensure that point and probabilistic forecast errors in Section 5 reflect genuine out-of-sample generalization from models fitted on 2010–2013 data and evaluated against held-out 2014 conditions, including the January 6–8 polar vortex event defined in Section 2.4.